% META: {"id":"1SPE-PRODSCAL-EV-A-corrige","chapitre":"1SPE-PRODUIT-SCALAIRE","type_objet":"corrige_evaluation","capacites_codes":["C1","C2","C3","C4","C5"],"status":"approved","evaluation_ref":"1SPE-PRODSCAL-EV-A","origin":"generated","status_history":["generated","audited","accepted","approved"],"release_acceptance":"RELEASE_OWNER_BATCH_ACCEPTANCE","acceptance_closure_digest":"sha256:9b3ccf9a81c5520fbb7e03b7b2d3e7bf2057a02834b6d908bf3be5f49f3b553f","acceptance_receipt_digest":"sha256:4fad86f0282e0fa4e0419cca1ad6379ae7af5a280c04a4a90880f90a1c044242"}

\section*{Corrige — Évaluation A — Produit scalaire}

%---------------------------------------------------------------
\subsection*{Exercice 1 — 5 points}
%---------------------------------------------------------------

\baremeIndicatif{Q1 : 1 pt ; Q2 : 1 pt ; Q3 : 1,5 pt ; Q4 : 1,5 pt}

\textbf{Question 1.}
$\vec{u} \cdot \vec{v} = 3 \times 2 + (-1) \times 5 = 6 - 5 = 1$.

\textbf{Question 2.}
$\|\vec{u}\| = \sqrt{9 + 1} = \sqrt{10}$ et $\|\vec{v}\| = \sqrt{4 + 25} = \sqrt{29}$.

\textbf{Question 3.}
$\|\vec{u} + \vec{v}\|^2 = \|\vec{u}\|^2 + 2\,\vec{u} \cdot \vec{v} + \|\vec{v}\|^2 = 10 + 2 + 29 = 41$.

Donc $\|\vec{u} + \vec{v}\| = \sqrt{41}$.

Vérification directe : $\vec{u} + \vec{v} = \begin{pmatrix} 5 \\ 4 \end{pmatrix}$, $\|\vec{u} + \vec{v}\|^2 = 25 + 16 = 41$. \checkmark

\textbf{Question 4.}
$(\vec{u} + \vec{v}) \cdot (\vec{u} - \vec{v}) = \|\vec{u}\|^2 - \|\vec{v}\|^2 = 10 - 29 = -19$.

Vérification : $\vec{u} + \vec{v} = \begin{pmatrix} 5 \\ 4 \end{pmatrix}$, $\vec{u} - \vec{v} = \begin{pmatrix} 1 \\ -6 \end{pmatrix}$. $(5)(1) + (4)(-6) = 5 - 24 = -19$. \checkmark

%---------------------------------------------------------------
\subsection*{Exercice 2 — 5 points}
%---------------------------------------------------------------

\baremeIndicatif{Q1 : 2 pts ; Q2 : 2 pts ; Q3 : 1 pt}

\textbf{Question 1.}
$\overrightarrow{AB}\begin{pmatrix} 3 \\ -2 \end{pmatrix}$ et $\overrightarrow{AC}\begin{pmatrix} 2 \\ 3 \end{pmatrix}$.

$\overrightarrow{AB} \cdot \overrightarrow{AC} = 3 \times 2 + (-2) \times 3 = 6 - 6 = 0$.

Le triangle $ABC$ est \textbf{rectangle en $A$}.

\textbf{Question 2.}
$\overrightarrow{BA}\begin{pmatrix} -3 \\ 2 \end{pmatrix}$, $\overrightarrow{BC}\begin{pmatrix} -1 \\ 5 \end{pmatrix}$.

$\overrightarrow{BA} \cdot \overrightarrow{BC} = 3 + 10 = 13$, $\|\overrightarrow{BA}\| = \sqrt{13}$, $\|\overrightarrow{BC}\| = \sqrt{26}$.

$\cos\widehat{ABC} = \dfrac{13}{\sqrt{13}\sqrt{26}} = \dfrac{13}{\sqrt{338}} = \dfrac{13}{13\sqrt{2}} = \dfrac{1}{\sqrt{2}} = \dfrac{\sqrt{2}}{2}$.

Donc $\widehat{ABC} = \dfrac{\pi}{4}$.

\textbf{Question 3.}
$\widehat{ACB} = \pi - \dfrac{\pi}{2} - \dfrac{\pi}{4} = \dfrac{\pi}{4}$.

Le triangle est rectangle isocele en $A$.

%---------------------------------------------------------------
\subsection*{Exercice 3 — 5 points}
%---------------------------------------------------------------

\baremeIndicatif{Q1 : 1,5 pt ; Q2 : 2 pts ; Q3 : 1,5 pt}

\textbf{Question 1.}
$I = (3; 0)$ milieu de $[AB]$ et $\overrightarrow{AB}\begin{pmatrix} 6 \\ 0 \end{pmatrix}$.

$\overrightarrow{IM} \cdot \overrightarrow{AB} = 0$ : $6(x - 3) + 0 = 0$, soit $x = 3$.

\textbf{Question 2.}
$\overrightarrow{BC}\begin{pmatrix} -4 \\ 4 \end{pmatrix}$. On parametrise : $H = B + t\,\overrightarrow{BC} = (6 - 4t; 4t)$.

$\overrightarrow{AH} \cdot \overrightarrow{BC} = 0$ : $(6-4t)(-4) + (4t)(4) = 0$, soit $-24 + 32t = 0$, d'ou $t = \dfrac{3}{4}$.

$H = (3; 3)$.

\textbf{Question 3.}
$\mathcal{A} = \dfrac{1}{2} |x_B \cdot y_C - x_C \cdot y_B| = \dfrac{1}{2} |6 \times 4 - 2 \times 0| = \dfrac{24}{2} = 12$.

%---------------------------------------------------------------
\subsection*{Exercice 4 — 5 points}
%---------------------------------------------------------------

\baremeIndicatif{Q1 : 2 pts ; Q2 : 1,5 pt ; Q3 : 1,5 pt}

\textbf{Question 1.}
$BC^2 = AB^2 + AC^2 - 2 \cdot AB \cdot AC \cdot \cos\widehat{BAC} = 64 + 25 - 80 \times \dfrac{1}{2} = 89 - 40 = 49$.

$BC = 7$.

\textbf{Question 2.}
$\cos\widehat{ABC} = \dfrac{BC^2 + AB^2 - AC^2}{2 \cdot BC \cdot AB} = \dfrac{49 + 64 - 25}{112} = \dfrac{88}{112} = \dfrac{11}{14}$.

\textbf{Question 3.}
$\mathcal{A} = \dfrac{1}{2} \times AB \times AC \times \sin\widehat{BAC} = \dfrac{1}{2} \times 8 \times 5 \times \dfrac{\sqrt{3}}{2} = 10\sqrt{3} \approx 17{,}3$.

% BEGIN-VERIFY
% from sympy import Matrix, sqrt, Rational, simplify, cos, sin, pi, symbols, solve
% # ---- Exercice 1 : u(3;-1), v(2;5) ----
% u = Matrix([3, -1]); v = Matrix([2, 5])
% assert u.dot(v) == 1
% assert simplify(u.norm() - sqrt(10)) == 0
% assert simplify(v.norm() - sqrt(29)) == 0
% assert (u + v).dot(u + v) == 41
% assert (u + v).dot(u - v) == -19
% assert (u + v).dot(u - v) == u.dot(u) - v.dot(v)
% # ---- Exercice 2 : A(1;2), B(4;0), C(3;5) ----
% A = Matrix([1, 2]); B = Matrix([4, 0]); C = Matrix([3, 5])
% AB = B - A; AC = C - A
% assert AB == Matrix([3, -2]) and AC == Matrix([2, 3])
% assert AB.dot(AC) == 0
% BA = A - B; BC = C - B
% assert BA.dot(BC) == 13
% assert simplify(BA.dot(BC)/(BA.norm()*BC.norm()) - sqrt(2)/2) == 0
% # la somme des angles ferme le triangle rectangle isocele
% assert simplify(pi - pi/2 - pi/4 - pi/4) == 0
% # ---- Exercice 3 : A(0;0), B(6;0), C(2;4) ----
% A = Matrix([0, 0]); B = Matrix([6, 0]); C = Matrix([2, 4])
% I = (A + B)/2
% assert I == Matrix([3, 0])
% t = symbols('t')
% H = B + t*(C - B)
% racine = solve((H - A).dot(C - B), t)
% assert racine == [Rational(3, 4)]
% H = H.subs(t, Rational(3, 4))
% assert H == Matrix([3, 3])
% aire = abs((B - A)[0]*(C - A)[1] - (C - A)[0]*(B - A)[1])/2
% assert aire == 12
% # ---- Exercice 4 : AB=8, AC=5, angle A = pi/3 ----
% AB_, AC_ = 8, 5
% BC2 = AB_**2 + AC_**2 - 2*AB_*AC_*cos(pi/3)
% assert simplify(BC2 - 49) == 0
% BC_ = 7
% assert simplify(Rational(BC_**2 + AB_**2 - AC_**2, 2*BC_*AB_) - Rational(11, 14)) == 0
% assert simplify(Rational(1, 2)*AB_*AC_*sin(pi/3) - 10*sqrt(3)) == 0
% END-VERIFY
